% Chapter 18: large and small. From Teddy Rook's chapter18.tex (November 2026), adapted to the paper's labels, macros and
% citation keys. Every printed number is recomputed by compute/c13_large_verify.py (his verify_ch18.py).
\chapter{Large and small: the same patterns in atoms and in orbits}\label{ch:large-small}

\addednote{\revD}

The Connections diagram of Chapter~\ref{ch:connections} has four subjects: music, patterns, mathematics and chemistry. This
chapter adds physics and astrophysics, and asks one question of them: where do the structures of the
earlier chapters appear both in very small things and in very large ones? We found four places where the
answer is exact or measurable, and a fifth section returns to Kepler, whose work runs through all of them.
None of the physics is ours. What we add is the reading through the continued fraction, several
computations, and two new exact statements about the correlation matrix
(Propositions~\ref{prop:18mirror} and~\ref{prop:18harmonic}).

A warning first. This subject is full of numerology: the Titius--Bode ``law'' of planetary distances and
Dirac's large-number coincidences are the classic cases of patterns that explain nothing. The discipline
of the earlier chapters applies. A pattern is reported with its test, and a coincidence is called one.

%-------------------------------------------------------------------------------------------------------
\section{The atom and the solar system are one problem}\label{sec:18kepler}

Section~\ref{sec:hydrogen} called the link between hydrogen's $1/n^{2}$ and $\zeta(2)$ ``a shared formula, not shared
physics,'' and Section~\ref{sec:shells} said the same of the shell sizes $2,8,18,32$. There is a neighbouring link
that \emph{is} shared physics, and it strengthens the chemistry corner of the diagram.

Gravity and the electric force both fall off as the inverse square of the distance. So a planet round the
Sun and, in Bohr's model, an electron round a proton obey the same law of periods.

\begin{proposition}[Kepler's third law, twice; classical]\label{prop:18kepler}
For a circular orbit of radius $r$ and period $T$,
\[
  T^{2}=\frac{4\pi^{2}}{GM}\,r^{3}\quad\text{(planet round a mass $M$)},\qquad
  T^{2}=\frac{4\pi^{2}m_e}{k e^{2}}\,r^{3}\quad\text{(electron round a proton)}.
\]
In Bohr's model the allowed orbits have $r_n=a_0n^{2}$ and therefore $T_n=T_1n^{3}$, with
$a_0=5.292\times10^{-11}$\,m and $T_1=1.520\times10^{-16}$\,s.
\end{proposition}

\begin{proof}
Set the attracting force equal to the centripetal force $m(2\pi/T)^{2}r$, with force $GMm/r^{2}$ in one
case and $ke^{2}/r^{2}$ in the other. Bohr's condition, angular momentum $n\hbar$, gives
$r_n=n^{2}\hbar^{2}/(m_eke^{2})$.
\end{proof}

Our script checks both. For the Earth, $T^{2}/r^{3}$ and $4\pi^{2}/(GM_\odot)$ are both
$2.9747\times10^{-19}$ in SI units; for the electron in its lowest orbit both sides are $0.15588$. The two
constants differ by a factor of $5\times10^{17}$, and the law is the same.

\paragraph{The rhythm of an atom, heard by a radio telescope.} An electron in a very high orbit behaves
almost classically, and the light it emits when it drops $k$ levels has $k$ times its orbital frequency:
the $k$-th harmonic of its rhythm. Astronomers have detected carbon atoms in interstellar space with
$n\approx1009$, through transitions with $k=4$ near $26$\,MHz \cite{stepkin2007}. The numbers fit. At
$n=1009$ the orbit has radius $a_0n^{2}=54\,\mu$m, so the atom is about a tenth of a millimetre across,
a million times the size of an ordinary atom; its orbital period is $T_1n^{3}=1.56\times10^{-7}$\,s, an
orbital frequency of $6.41$\,MHz; and the transition $n=1013\to1009$ has frequency $25.5$\,MHz, which is
$4.000$ times the orbital frequency at $n=1011$. A dish on Earth is listening to the fourth harmonic of
an electron's orbit.

\paragraph{The hidden symmetry, and the periodic table.} An inverse-square force has a conserved quantity
beyond energy and angular momentum, the Laplace--Runge--Lenz vector, which points along the long axis of
the ellipse. It is the reason planetary orbits are closed ellipses that do not precess. Pauli (1926) and
Fock (1935) showed that the same conserved vector gives the hydrogen atom a four-dimensional rotational
symmetry, and that this symmetry is why all $n^{2}$ states of level $n$ have the same energy
\cite{pauli1926,fock1935}. With two spin states each, that is $2n^{2}$: the shell sizes $2,8,18,32$. So
the squares of Section~\ref{sec:shells} come from the symmetry that keeps the planets on closed ellipses. In atoms
with many electrons the other electrons screen the nucleus, the force is no longer exactly inverse-square,
and the symmetry is broken; the rows of the periodic table then have lengths $2,8,8,18,18,32,32$, each
shell size after the first appearing twice. The periodic table is hydrogen's symmetry, broken.

\paragraph{A historical echo, and a caution.} Kepler published his third law in \emph{Harmonices Mundi}
(1619), a book whose subject is the music of the planets: he assigned each planet a melody from the ratio
of its fastest and slowest speeds. And the exponent in $T\propto r^{3/2}$ is the ratio of a musical
fifth. That last fact is a coincidence of small numbers and nothing more.

%-------------------------------------------------------------------------------------------------------
\section{Resonance as rhythm}\label{sec:18resonance}

Section~\ref{sec:locking} treated mode locking: coupled oscillators fall into step at simple frequency ratios. Orbits do
the same. Two bodies whose periods are near a ratio $p/q$ meet in the same configuration every $q$ turns
of one and $p$ of the other, and the repeated tug can lock the ratio. The table of convergents of Chapter
2 turns out to be a detector for this.

\paragraph{The big-term test.} Expand the ratio of two periods as a continued fraction. If it is locked
to $p/q$, then $p/q$ is a convergent and \emph{the next term is large}, for the reason that $292$ follows
$355/113$ in $\pi$. Table~\ref{tab:18res} applies this to four systems. The planets' periods are their long-term mean
periods \cite{standish1992}, the moons' are from NASA's fact sheets \cite{nasafact}, and the others are
from the discovery papers \cite{agol2021,luque2023}.

\begin{table}[ht]
\centering
\caption{Neighbouring period ratios (outer over inner): the expansion, the nearest fraction with
denominator at most $8$, the normalized error $q^{2}|r-p/q|$ of \cref{tab:pi-conv}, and the term that follows
$p/q$. The last line of each block compares the mean normalized error with that of random ratios drawn
log-uniformly from $[1.25,3]$.}\label{tab:18res}
\small
\begin{tabular}{llllrr}
\hline
pair & ratio & expansion & nearest & $q^{2}\times$error & next term\\
\hline
Venus/Mercury & 2.5543 & $[2;1,1,4,\ldots]$ & 5/2 & 0.217 & 4\\
Earth/Venus & 1.6255 & $[1;1,1,1,2,\ldots]$ & 13/8 & 0.033 & 29\\
Mars/Earth & 1.8808 & $[1;1,7,2,\ldots]$ & 2/1 & 0.119 & 7\\
Jupiter/Mars & 6.3067 & $[6;3,3,1,\ldots]$ & 19/3 & 0.239 & 3\\
Saturn/Jupiter & 2.4833 & $[2;2,14,2,\ldots]$ & 5/2 & 0.067 & 14\\
Uranus/Saturn & 2.8521 & $[2;1,5,1,\ldots]$ & 3/1 & 0.148 & 5\\
Neptune/Uranus & 1.9614 & $[1;1,24,1,\ldots]$ & 2/1 & 0.039 & 24\\
Pluto/Neptune & 1.5046 & $[1;1,1,53,\ldots]$ & 3/2 & 0.018 & 53\\
\multicolumn{6}{l}{\emph{Solar system: mean $0.110$; random $0.146$; chance of being this close $0.14$}}\\
\hline
Europa/Io & 2.0073 & $[2;137,\ldots]$ & 2/1 & 0.007 & 137\\
Ganymede/Europa & 2.0147 & $[2;68,\ldots]$ & 2/1 & 0.015 & 68\\
Callisto/Ganymede & 2.3326 & $[2;3,160,\ldots]$ & 7/3 & 0.006 & 160\\
\multicolumn{6}{l}{\emph{Jupiter's moons: mean $0.009$; chance $0.0001$}}\\
\hline
TRAPPIST-1 c/b & 1.6031 & $[1;1,1,1,\ldots]$ & 8/5 & 0.076 & 12\\
d/c & 1.6719 & $[1;1,2,20,\ldots]$ & 5/3 & 0.047 & 20\\
e/d & 1.5067 & $[1;1,1,36,\ldots]$ & 3/2 & 0.027 & 36\\
f/e & 1.5092 & $[1;1,1,26,\ldots]$ & 3/2 & 0.037 & 26\\
g/f & 1.3416 & $[1;2,1,12,\ldots]$ & 4/3 & 0.074 & 12\\
h/g & 1.5198 & $[1;1,1,12,\ldots]$ & 3/2 & 0.079 & 12\\
\multicolumn{6}{l}{\emph{TRAPPIST-1: mean $0.057$; chance $0.005$}}\\
\hline
HD 110067 c/b & 1.5003 & $[1;1,1,\ldots]$ & 3/2 & 0.001 & $>300$\\
d/c & 1.5007 & $[1;1,1,\ldots]$ & 3/2 & 0.003 & $>300$\\
e/d & 1.5007 & $[1;1,1,\ldots]$ & 3/2 & 0.003 & $>300$\\
f/e & 1.3334 & $[1;2,1,\ldots]$ & 4/3 & 0.000 & $>300$\\
g/f & 1.3339 & $[1;2,1,\ldots]$ & 4/3 & 0.006 & 180\\
\multicolumn{6}{l}{\emph{HD 110067: mean $0.0025$; chance below $0.0001$}}\\
\hline
\end{tabular}
\end{table}

\paragraph{What the table shows.} The true resonances announce themselves with large terms: $53$ for
Pluto and Neptune (locked at $3/2$), $137$, $68$ and $160$ for Jupiter's moons, and from $180$ to over
$3000$ for HD~110067. (Of the three for Jupiter, $137$ and $68$
belong to the locked trio Io, Europa and Ganymede. Callisto's $160$ marks a very near miss of $7/3$ that
is not known to be a dynamical lock, a reminder that a large term is evidence and not proof.) Three pairs of Saturn's moons
do the same: Tethys/Mimas $[2;318,\ldots]$, Dione/Enceladus $[1;1,388,\ldots]$ and Hyperion/Titan
$[1;2,1,109,\ldots]$. The planets of the solar system taken together are \emph{not} significantly closer
to simple fractions than random ratios would be ($p=0.14$); the chains of TRAPPIST-1 and HD~110067 are,
beyond doubt. Two near-misses in the solar system are famous for their consequences and not locked:
Saturn/Jupiter, $[2;2,14,\ldots]$, near $5/2$, drives a swing in the motions of the two planets with a
period of $883$ years, known as the great inequality, and
Neptune/Uranus sits near $2/1$. Earth/Venus is $[1;1,1,1,2,\ldots]$, the start of the golden ratio's
expansion, with convergents $8/5$ and $13/8$: the reason Venus traces a five-petalled figure in our sky
every eight years. It is not a resonance in the dynamical sense. (The $137$ of Io and Europa has nothing
to do with the $137$ of atomic physics.)

\paragraph{Three bodies in step.} For Io, Europa and Ganymede the mean motions (degrees per day) satisfy
$n_1-3n_2+2n_3=-0.00003$, zero to the accuracy of the data, while the two-body near-misses $n_1-2n_2$ and
$n_2-2n_3$ are both $0.7395$: the two pairs drift off the exact ratio $2/1$ at the same rate, so the
three-body pattern repeats exactly. This is Laplace's resonance. As a polyrhythm it is $4$ against $2$
against $1$ with a slow common drift.

\paragraph{As music.} The resonant chains are stacked intervals in the sense of Section~\ref{sec:tuning}. HD~110067 is
three fifths and two fourths ($3/2,3/2,3/2,4/3,4/3$): while the outermost planet goes round $9$ times the
others go round $12$, $16$, $24$, $36$ and $54$ times. TRAPPIST-1 is looser: while its outermost planet
goes round twice the others make $3.04$, $4.08$, $6.15$, $9.27$, $15.5$ and $24.9$ turns, close to
$3:4:6:9:15:24$. Either can be loaded into the polyrhythm instrument of the companion site as it stands.

\paragraph{Where resonance removes instead of locks.} Kepler's third law puts the orbits whose periods
are $1/3$, $2/5$, $3/7$ and $1/2$ of Jupiter's at $2.500$, $2.824$, $2.957$ and $3.277$ astronomical
units. Those are the Kirkwood gaps of the asteroid belt, where almost no asteroids are found. (A commonly
quoted set of positions is $2.502$, $2.825$, $2.958$ and $3.279$. It comes from scaling Jupiter's distance
instead of using its period, which adds about $0.0015$ units: Jupiter's own mass enters its orbit but not
the asteroid's, and the value taken for Jupiter's distance varies. A resonance is a statement about
periods, so we give the period-based figures.) In the same way an orbit with half the period of Saturn's
moon Mimas lies, by Kepler's law, $116{,}780$\,km from Saturn's centre. The bright B ring ends and the
Cassini Division begins at $117{,}580$\,km, $0.7$ per cent further out. The difference is Saturn's
flattened shape, which Kepler's law ignores: with the measured flattening put into the calculation the
resonance moves to $117{,}530$\,km, within $50$\,km of the edge. So simple fractions sometimes lock an orbit
(Pluto, the moons) and sometimes empty it (the gaps). The theorem behind the second case is that of
Kolmogorov, Arnold and Moser: an orbit survives small disturbances when its frequency ratio is hard to
approximate by fractions, that is, when its continued fraction has small terms. The golden ratio, all of
whose terms are $1$, is the hardest number of all to approximate (Chapter~\ref{ch:cf}); a number with a $292$ in it
is easy to approximate by the fraction that precedes the $292$.

%-------------------------------------------------------------------------------------------------------
\section{The mirror in the correlation matrix}\label{sec:18mirror}

\Cref{thm:th-top,thm:th-min} say that the largest eigenvalue of the correlation matrix $C_m$ and the reciprocal
of the smallest grow at the same rate. Is there an exact symmetry behind that? There is, in two parts.

\begin{proposition}[the sign mirror]\label{prop:18mirror}
Let $\lambda(a)=(-1)^{\Omega(a)}$ be Liouville's function, where $\Omega(a)$ counts the prime factors of
$a$ with multiplicity. For $2\le a,b\le m$,
\[
  \bigl(C_m^{-1}\bigr)_{ab}=\lambda(a)\lambda(b)\,\sqrt{(a-1)(b-1)}
     \sum_{\substack{d\le m,\ \operatorname{lcm}(a,b)\mid d\\ d/a,\ d/b\ \text{square-free}}}\frac{1}{\varphi(d)} .
\]
Hence, with $S=\operatorname{diag}(\lambda(a))$, every entry of $SC_m^{-1}S$ is non-negative; $1/\lambda_{\min}(C_m)$
is the Perron eigenvalue of that non-negative matrix; and an eigenvector of $C_m$ for its smallest
eigenvalue can be chosen of the form $\lambda(a)w_a$ with every $w_a\ge0$.
\end{proposition}

\begin{proof}
In the formula of \cref{prop:th-inv} a term is non-zero only when $d/a$ and $d/b$ are square-free, and then
$\mu(d/a)=\lambda(d/a)=\lambda(d)\lambda(a)$, because $\lambda$ is completely multiplicative and
$\lambda(a)^{2}=1$. So $\mu(d/a)\mu(d/b)=\lambda(a)\lambda(b)$ for every term. The rest is the
Perron--Frobenius theorem.
\end{proof}

So the two ends of the spectrum are tied by a change of sign. The top eigenvector has no negative entries: all the
pulse trains in phase. The bottom one is a positive vector with the sign $\lambda(a)$: every pulse train
flipped once for each prime factor. And the two matrices are built in mirrored ways. An entry of $C_m$ is
a sum of $\varphi(d)$ over the \emph{common divisors} $d$ of $a$ and $b$ (\cref{prop:oq-gram}); an entry of
the inverse is, up to sign, a sum of $1/\varphi(d)$ over their \emph{common multiples} up to $m$. What two
rhythms share below, against where they meet again above: $\gcd$ against $\operatorname{lcm}$.

\begin{proposition}[one prime at a time; classical]\label{prop:18chain}
Let $G=[\gcd(a,b)/\sqrt{ab}]$, the matrix of Chapter~\ref{ch:third} that differs from $C_m$ only in the shift
from $\gcd$ to $\gcd-1$. Restricted to the powers $1,p,p^{2},\ldots$ of one prime $p$, it is $K=[r^{|i-j|}]$ with
$r=p^{-1/2}$, the correlation matrix of a chain with correlation $r$ between neighbours. Its inverse is
tridiagonal. On the infinite chain its spectrum is the interval $[1/M,M]$ with $M=(\sqrt p+1)/(\sqrt p-1)$,
and the map
\[
  \lambda\;\longmapsto\;\frac{\lambda}{2c\lambda-1},\qquad c=\frac{p+1}{p-1},
\]
carries the spectrum to itself and exchanges its two ends.
\end{proposition}

\begin{proof}
On the infinite chain $K$ is convolution with $r^{|k|}$, whose Fourier series is the Poisson kernel
$P(\theta)=(1-r^{2})/(1-2r\cos\theta+r^{2})$, with range $[(1-r)/(1+r),(1+r)/(1-r)]$. Then
$1/P(\theta)+1/P(\theta+\pi)=2(1+r^{2})/(1-r^{2})=2c$, which is the stated map, and $1/P$ has only the
frequencies $0,\pm1$, which is the tridiagonal inverse.
\end{proof}

The shift $\theta\mapsto\theta+\pi$ is multiplication by $(-1)^{k}$, Liouville's sign on the powers of
$p$. So prime by prime the mirror is exact, and it is a fractional-linear map, the same kind of map that
drives continued fractions. Because $\gcd(a,b)/\sqrt{ab}$ is multiplicative, the full matrix $G$ is the product of these chains over
all primes, cut off at $m$. Without the cut-off the top would be $\prod_pM_p=\infty$ and the bottom $0$;
the cut-off makes both finite, and \cref{thm:th-top,thm:th-min} say that for $C_m$ it treats them alike to
leading order. It does not treat them exactly alike: $1/\lambda_{\min}$ is about $2.6$ times $\lambda_{\max}$
(Table~\ref{tab:18ratio} below), and the mean of the eigenvalues is $1$ while the mean of their reciprocals
tends to $3.28$ (\cref{prop:th-trace}).

\paragraph{The numbers.} At $m=400$ the smallest entry of $SC_m^{-1}S$ is $-9\times10^{-15}$, zero to
rounding, and the formula of Proposition~\ref{prop:18mirror} matches a dense inverse to $10^{-14}$. At
$m=20{,}000$ the bottom eigenvector times $\lambda(a)$ is non-negative on every index. The two vectors
live in different places: the top one is largest at $a=60,120,180,24,12,\ldots$, numbers with many
repeated small factors, and the bottom one at $a=30,6,60,42,12,210,\ldots$, led by the products of the
first few distinct primes.

\paragraph{In physics.} The chain of Proposition~\ref{prop:18chain} is the correlation matrix of the
one-dimensional Ising model of magnetism, and of a damped random walk; a tridiagonal inverse is the
statement that each site depends only on its neighbours. So the arithmetic correlation matrix is a
lattice of such chains, one direction for each prime, and its inverse is local: the entry for $a$ and $b$
vanishes unless they have a common multiple $d\le m$ with $d/a$ and $d/b$ square-free. Physics has exact
small-to-large symmetries of this kind, the best known being the duality that exchanges high and low
temperature in the two-dimensional Ising model. We claim only the analogy.

\paragraph{Both ends, written with common beats.} Proposition~\ref{prop:18mirror} wrote the bottom of the
spectrum as a sum over common multiples. The top can be written the same way.

\begin{proposition}[the harmonic form]\label{prop:18harmonic}
$C_m$ has the same eigenvalues as the non-negative matrix
\[
  T_m=\Bigl[\sqrt{\varphi(a)\varphi(b)}\sum_{\substack{n\le m\\ \operatorname{lcm}(a,b)\mid n}}\frac{1}{n-1}\Bigr]_{2\le a,b\le m},
\]
and $G_m=[\gcd(a,b)/\sqrt{ab}]_{1\le a,b\le m}$ has the same eigenvalues as
\[
  \Bigl[\frac{\sqrt{\varphi(a)\varphi(b)}}{\operatorname{lcm}(a,b)}\;H_{\lfloor m/\operatorname{lcm}(a,b)\rfloor}\Bigr]_{1\le a,b\le m},
  \qquad H_k=1+\tfrac12+\cdots+\tfrac1k .
\]
\end{proposition}

\begin{proof}
By \cref{prop:oq-gram}, $C_m=DV\Phi V^{T}D=B^{T}B$ with $B=\Phi^{1/2}V^{T}D$, a square matrix. So $C_m$ and
$BB^{T}$ have the same eigenvalues, and the $(a,b)$ entry of $BB^{T}$ is
$\sqrt{\varphi(a)\varphi(b)}\sum_{n}[a\mid n][b\mid n]/(n-1)$. For $G_m$ the weight $1/(n-1)$ becomes $1/n$,
and the sum of $1/n$ over the multiples of $L$ up to $m$ is $H_{\lfloor m/L\rfloor}/L$.
\end{proof}

The number $\operatorname{lcm}(a,b)$ is the beat on which the pulse trains $a$ and $b$ first coincide, the
length of their polyrhythm (Section~\ref{sec:polyrhythm}), and its multiples are their later coincidences. So both ends of
the spectrum are Perron eigenvalues of matrices that count coincidences up to the horizon $m$. For the top
every coincidence $n$ counts, with weight $1/(n-1)$. For the bottom
(Proposition~\ref{prop:18mirror}) the weight is $1/\varphi(n)$, and a coincidence counts only if $n/a$ and
$n/b$ are square-free. In the gcd matrix the first of these is the harmonic series itself: the entry for
two rhythms is a harmonic number, $H_k$, where $k$ is how many times they coincide before the horizon.

\paragraph{How far apart are the two ends?} With these forms the two ends can be computed for much larger
matrices (the divisor sums are done one prime at a time). Table~\ref{tab:18ratio} gives the result.

\begin{table}[ht]
\centering
\caption{The two ends of the spectrum and their ratio $(1/\lambda_{\min})/\lambda_{\max}$, for the
correlation matrix $C_m$ and, in the last column, for the gcd matrix $G_m$.}\label{tab:18ratio}
\small
\begin{tabular}{rrrrr}
\hline
$m$ & $\lambda_{\max}(C_m)$ & $1/\lambda_{\min}(C_m)$ & ratio for $C_m$ & ratio for $G_m$\\
\hline
$10^{2}$ & 8.096 & 21.967 & 2.713 & 2.515\\
$10^{3}$ & 18.773 & 50.976 & 2.715 & 2.550\\
$10^{4}$ & 36.429 & 97.857 & 2.686 & 2.570\\
$10^{5}$ & 63.078 & 167.104 & 2.649 & 2.572\\
$10^{6}$ & 101.252 & 265.276 & 2.620 & 2.570\\
$10^{7}$ & 154.040 & 400.477 & 2.600 & 2.568\\
$3\times10^{7}$ & 185.434 & 480.554 & 2.592 & 2.565\\
\hline
\end{tabular}
\end{table}

Three things stand out. First, for the gcd matrix every value we computed from $m=2\times10^{4}$ to $3\times10^{7}$, some fifty
sizes in all, lies between $2.565$ and $2.573$, while both eigenvalues grew fourfold. Second, the ratio for $C_m$
is falling towards the ratio for $G_m$: the difference between them is $0.165$, $0.117$, $0.077$, $0.050$,
$0.032$ and $0.026$ down the table from $10^{3}$. So the question is really one about $G_m$. Third,
neither ratio has settled: since $10^{5}$ the ratio for $G_m$ has been drifting down by about $0.002$ for
each tenfold increase in $m$.

\begin{proposition}[two cases where the ratio is exact]\label{prop:18exact}
(i) Let $n$ be square-free. On the divisors of $n$ the gcd matrix has
$\lambda_{\max}=\prod_{p\mid n}(1+p^{-1/2})$ and $\lambda_{\min}=\prod_{p\mid n}(1-p^{-1/2})$, so
\[
  \frac{1/\lambda_{\min}}{\lambda_{\max}}=\prod_{p\mid n}\frac{p}{p-1}=\frac{n}{\varphi(n)} .
\]
(ii) On the powers $1,p,\ldots,p^{k-1}$ of one prime the ratio is $p/(p-1)$ for $k=2$ and for $k=3$, and it
tends to $1$ as $k\to\infty$.
\end{proposition}

\begin{proof}
(i) The matrix is the Kronecker product of the $2\times2$ matrices with rows $(1,r)$ and $(r,1)$,
$r=p^{-1/2}$, whose eigenvalues are $1\pm r$. (ii) For $k=3$ the eigenvalues of
Proposition~\ref{prop:18chain}'s matrix are $1-r^{2}$ and $\tfrac12\bigl(2+r^{2}\pm r\sqrt{r^{2}+8}\bigr)$; the
last two are the extreme ones and their product is $1-r^{2}$. As $k\to\infty$ the extremes tend to $M$ and
$1/M$.
\end{proof}

So the ratio measures how short the chains are. Where each prime is used once, it is the totient ratio
$n/\varphi(n)$ of Chapter~\ref{ch:coprime}; where a prime is used to a high power, it falls back to $1$. Computing with
only the numbers built from the primes up to a bound $y$ shows how the two effects combine
(Table~\ref{tab:18smooth}).

\begin{table}[ht]
\centering
\caption{The gcd matrix on the numbers up to $10^{7}$ whose prime factors are all at most $y$. With a
fixed set of primes both ends have the same ceiling, $\prod_{p\le y}M_p$.}\label{tab:18smooth}
\small
\begin{tabular}{rrrrrr}
\hline
$y$ & numbers & $\lambda_{\max}$ & $1/\lambda_{\min}$ & ratio & ceiling\\
\hline
2 & 24 & 5.335 & 5.805 & 1.088 & 5.83\\
3 & 190 & 15.735 & 21.056 & 1.338 & 21.8\\
5 & 768 & 29.841 & 50.513 & 1.693 & 56.9\\
7 & 2\,155 & 45.430 & 94.498 & 2.080 & 126.2\\
13 & 8\,289 & 69.097 & 175.848 & 2.545 & 415.5\\
29 & 37\,627 & 96.862 & 261.104 & 2.696 & 2\,418\\
97 & 269\,882 & 127.264 & 338.751 & 2.662 & $1.3\times10^{5}$\\
997 & 2\,028\,358 & 150.352 & 391.162 & 2.602 & \\
all & $10^{7}$ & 157.790 & 405.157 & 2.568 & $\infty$\\
\hline
\end{tabular}
\end{table}

The table shows the mechanism. \emph{The bottom reaches its ceiling sooner than the top.} With the powers
of $2$ alone, $24$ levels bring the bottom to within half a per cent of the ceiling $3+2\sqrt2=5.83$,
while the top is still more than eight per cent short. Each further small prime adds to the bottom's lead, up to
about $2.7$ near $y=29$. The large primes then take a little of it back, because they still add to the
top while the bottom has less left to gain. The same picture shows in the eigenvectors. At $m=10^{6}$
half the weight of the top vector lies on indices above $47{,}000$, and half the weight of the bottom
vector on indices below $600$: the bottom is made by a few small rhythms, the top by many large ones.

For any fixed set of primes the ratio falls to $1$ as $m$ grows: with powers of $2$ only it is $1.200$,
$1.117$ and $1.088$ at $m=10^{4}$, $10^{6}$ and $10^{7}$. With all primes it stays near $2.57$, because new
primes keep arriving as the old ones saturate. Whether that balance has a limit is the open question, and
the computation cannot decide it. A limit near $2.5$ and a very slow decline are both consistent with
Table~\ref{tab:18ratio}. We did not recognize $2.57$ as any known constant, and at these sizes it would be
numerology to try.

%-------------------------------------------------------------------------------------------------------
\section{The butterfly of a number}\label{sec:18butterfly}

An electron moving on a square crystal lattice in a magnetic field obeys Harper's equation
\[
  \psi_{k+1}+\psi_{k-1}+2\cos(2\pi k\alpha+\nu)\,\psi_k=E\,\psi_k ,
\]
where $\alpha$ is the magnetic flux through one cell of the lattice in natural units. Hofstadter (1976)
found that when $\alpha=p/q$ the allowed energies form $q$ bands, and that the picture over all $\alpha$
is a fractal, the ``butterfly'' \cite{hofstadter1976}. It was observed in graphene superlattices in 2013
\cite{dean2013}. For irrational $\alpha$ the spectrum is a Cantor set \cite{avila2009}, like the dust of
Section~\ref{sec:dust}.

The spectrum of $\alpha$ is approached through the convergents of $\alpha$, so the continued fraction
organizes it. We computed the bands at the convergents of the five numbers of the Suite
(Table~\ref{tab:18bands}).

\begin{table}[ht]
\centering
\caption{Energy bands of Harper's equation at convergents $p/q$: the sizes of the groups into which the
widest gaps divide the $q$ bands, and the total width of the bands against Thouless's estimate
$32G/(\pi q)=9.33/q$ ($G$ is Catalan's constant).}\label{tab:18bands}
\small
\begin{tabular}{llllr}
\hline
number & $p/q$ & groups & total width & $9.33/q$\\
\hline
$\pi-3$ & 1/7 & $1,1,1,1,1,1,1$ & 1.2845 & 1.3329\\
 & 15/106 & $15,15,15,16,15,15,15$ & 0.0888 & 0.0880\\
 & 16/113 & $16,16,16,17,16,16,16$ & 0.0825 & 0.0826\\
$\varphi-1$ & 8/13 & $5,3,5$ & 0.7116 & 0.7177\\
 & 13/21 & $8,5,8$ & 0.4438 & 0.4443\\
 & 21/34 & $13,8,13$ & 0.2753 & 0.2744\\
$\sqrt2-1$ & 12/29 & $5,7,5,7,5$ & 0.3216 & 0.3217\\
 & 29/70 & $12,17,12,17,12$ & 0.1333 & 0.1333\\
$e-2$ & 51/71 & $20,31,20$ & 0.1314 & 0.1314\\
$\log_2(3/2)$ & 7/12 & $2,3,2,3,2$ & 0.7832 & 0.7775\\
 & 31/53 & $9,13,9,13,9$ & 0.1760 & 0.1760\\
\hline
\end{tabular}
\end{table}

For $\pi$ the $113$ bands fall into seven groups of $16$, with one of $17$ in the middle: the bar of $7$
and the phrase of $113$ of Chapter~\ref{ch:rhythm}, as an energy spectrum. For the golden ratio the groups are
Fibonacci numbers, for $\sqrt2$ they are built from the Pell numbers $2,5,12,29$ (with $7=5+2$ and $17=12+5$), and for the fifth the twelve bands of
$7/12$ divide as $2,3,2,3,2$. The total width shrinks like $9.33/q$ \cite{thouless1983}, which is the
Cantor dust appearing.

Each gap carries a whole number. If $r$ bands lie below a gap, there is exactly one pair of integers
$(s,t)$ with $|t|\le q/2$ and
\[
  r=q\,s+p\,t ,
\]
found by Euclid's algorithm, and $t$ is the Hall conductance of the electron gas when the gap is at the
Fermi energy, in units of $e^{2}/h$ \cite{tknn1982}. For the six widest gaps of $16/113$ these integers
are $1,2,3,-3,-2,-1$. So the oldest algorithm in the book produces a quantity that laboratories measure
to parts in a billion.

An interactive version, in which any continued fraction can be entered, its levels marked on the
butterfly and its bands played from low energy to high, accompanies this chapter.

%-------------------------------------------------------------------------------------------------------
\section{Kepler, continued}\label{sec:18keplercont}

Kepler has now appeared twice: his third law in Section~\ref{sec:18kepler} and his harmony of the world
in Section~\ref{sec:18resonance}. Three more pieces of his work meet the methods of this paper, and in
each the continued fraction adds something he did not have.

\paragraph{The trigon is a continued fraction.} On 19 July 1595 Kepler drew for his class in Graz the places
in the zodiac where Jupiter and Saturn meet, about every twenty years \cite{gingerich1973}. Each meeting falls roughly two thirds of
the way round from the last, so three of them make a triangle, the \emph{trigon}, and the triangle slowly
turns. The drawing gave him the idea for his first book \cite{kepler1596}. It is a picture of a continued
fraction.

\begin{proposition}[conjunction figures]\label{prop:18conj}
Let two bodies move on circular orbits in one plane with periods $T_1<T_2$, and let $r=T_2/T_1$.
Successive conjunctions are separated by the time $T_1T_2/(T_2-T_1)$, and each lies
\[
  \alpha=\frac{1}{r-1}
\]
of a turn beyond the last. Hence:\\
(i) the first $N$ conjunctions cut the circle into arcs of at most three different lengths, for every
$N$;\\
(ii) if $p/q$ is a convergent of the fractional part of $\alpha$ and $a$ is the term that follows it, then
every $q$ conjunctions nearly close a figure with $q$ points, which has turned by $q\alpha-p$ of a circle;
and the figure has turned through the spacing of its own points, $1/q$ of a circle, after
$1/(q^{2}|\alpha-p/q|)$ rounds, a number between $a$ and $a+2$.
\end{proposition}

\begin{proof}
Between conjunctions the inner body gains one turn on the outer, so the time $t$ satisfies
$t/T_1-t/T_2=1$, and in that time the outer body moves $t/T_2=1/(r-1)$ turns. The conjunction points are
therefore the orbit of a rotation by $\alpha$. (i) is the three-gap theorem of Section~\ref{sec:three-gaps}, and (ii) is
the bound $1/(q(q'+q))<|\alpha-p/q|<1/(qq')$ of Chapter~\ref{ch:cf}, where $q'=aq+q''$ is the next denominator and
$q''\le q$ the previous one.
\end{proof}

So the figure is the convergent, and \emph{the next term counts how many times the figure is drawn before
it has turned by one point}. This is the normalized error of \cref{tab:pi-conv} again. Table~\ref{tab:18conj}
applies it to the planets, with long-term mean periods; real orbits are not circles, so real conjunctions wander a
few degrees about these places.

\begin{table}[ht]
\centering
\caption{Conjunction figures from mean periods. Each line gives the step $\alpha$ between meetings, its
expansion, and the cycles that the convergents produce. An opposition of an outer planet is a conjunction
of the Earth with it, seen from the Sun.}\label{tab:18conj}
\small
\begin{tabular}{llp{0.47\textwidth}}
\hline
pair & step $\alpha$ & cycles\\
\hline
Jupiter, Saturn & $0.6742=[0;1,2,14,2,\ldots]$ & $2/3$: a triangle every $59.6$ years, turned by
$8.1^{\circ}$; by $120^{\circ}$ after $14.8$ rounds, $883$ years. $29/43$: $853.9$ years.\\
Venus, Earth & $0.5987=[0;1,1,2,29,2,\ldots]$ & $3/5$: a five-pointed star every $7.99$ years, turned
by $2.4^{\circ}$; by $72^{\circ}$ after $29.9$ rounds. $88/147$: $235.0$ years. $91/152$: $243.00$ years.\\
Earth, Mars & $0.1353=[0;7,2,1,1,3,\ldots]$ & $7$ oppositions in $14.95$ years; $15$ in $32.0$; $22$ in
$47.0$; $37$ in $79.01$.\\
Earth, Jupiter & $0.0921=[0;10,1,6,4,\ldots]$ & $11$ oppositions in $12.01$ years; $76$ in $83.00$.\\
Earth, Saturn & $0.0351=[0;28,2,5,\ldots]$ & $28$ oppositions in $29.0$ years; $57$ in $59.00$.\\
Mercury, Earth & $0.3172=[0;3,6,1,1,2,1,\ldots]$ & $22$ meetings in $6.98$ years; $41$ in $13.01$; $104$
in $32.99$; $145$ in $46.00$.\\
\hline
\end{tabular}
\end{table}

The first line is Kepler's trigon. The convergent $2/3$ is the triangle; the term $14$ that follows says
the triangle is drawn about fifteen times before it has turned onto itself, which takes $883$ years. That
is exactly the period of the great inequality of Section~\ref{sec:18resonance}: the turning of the trigon
and the slow swing of Jupiter and Saturn are one beat, with frequency $5n_S-2n_J$ in the two mean
motions, seen in two ways. [Note: the figure turns in $883$ years; the swing it drives runs a few per cent slower, because the long axes of the two orbits move. See Section~\ref{sec:nw-sim}. \addedon{\revF}] This is the ``great cycle'' of the conjunctions, which medieval astrologers
put at $960$ years and Kepler at $794$. Kepler's figure is forty conjunctions ($794.4$ years), and $27/40$
is an intermediate fraction, the last step before the next convergent $29/43$. The term $14$ is delicate.
Jupiter disturbs Saturn's motion, and with the period of Saturn's present orbit in place of its long-term
mean the term is $13$ and the turn takes $842$ years. The second line is the five-pointed star of Venus,
with the large term $29$ behind it; the fraction $91/152$, the first intermediate step beyond the
convergent $88/147$, gives $243.00$ years, the period after which transits of Venus across
the Sun repeat. Kepler predicted the transit of Venus of 1631. The remaining lines
reproduce periods that Babylonian astronomers already used to predict the planets: $8$ years for Venus,
$46$ for Mercury, $47$ and $79$ for Mars, $83$ for Jupiter and $59$ for Saturn \cite{gray2009}.
Their second period for Jupiter, $71$ years, is $65$ oppositions: the fraction $6/65$ is not a convergent
but an intermediate fraction, the fifth of the six steps that lead from $1/10$ to $7/76$, on one of the
veins of Section~\ref{sec:th6}.

The figures themselves are Kepler's too. In the first two books of \emph{Harmonices Mundi} he studied the
regular polygons, star polygons among them, and how they fit together; the figure for a convergent $p/q$
is the star polygon with $q$ points joined $p$ apart.
He had both halves, the stars and the conjunctions; the continued fraction is what joins them. For
Jupiter's moons the same proposition gives the resonance of Section~\ref{sec:18resonance}: the meeting
points of Io with Europa, and of Europa with Ganymede, step backwards round Jupiter at the same rate,
$0.7395^{\circ}$ a day, once round in $487$ days.

An interactive version, ``Kepler's Figures'', accompanies this section: any two bodies can be chosen and
their meetings watched as they draw the figure and turn it. With Kepler's ellipses in place of circles it
puts the meeting of Venus and the Earth at $2012.43$, the transit of June 2012, and the meeting of Jupiter
and Saturn as seen from the Sun at $2020.84$, seven weeks before the two were seen together from the
Earth.

\paragraph{Kepler's harmonies, tested.} In Book V of \emph{Harmonices Mundi} \cite{kepler1619} Kepler
compared each planet's fastest angular speed as seen from the Sun, at perihelion, with its slowest, at
aphelion, and heard a musical interval in each ratio. For an orbit of eccentricity $e$ the ratio is
$((1+e)/(1-e))^{2}$. He printed the speeds he used (for Saturn $1'46''$ and $2'15''$ of arc a day, for the
Earth $57'3''$ and $61'18''$), and he printed his own misses beside them: Saturn's ratio exceeds a major
third by $53:54$, Mars falls short of a fifth by $29:30$. Table~\ref{tab:18harm} sets three versions side
by side: his own figures, the orbits as they were in 1600, and the orbits of today \cite{standish1992}.

\begin{table}[ht]
\centering
\caption{Kepler's intervals for the single planets, and by how much the ratio of fastest to slowest
angular speed misses each one, in cents (a semitone of the piano is $100$ cents). The orbits of 1600 and
2000 are mean orbits from one formula valid from 3000 BC to 3000 AD.}\label{tab:18harm}
\small
\begin{tabular}{llrrrr}
\hline
planet & Kepler's interval & his figures & 1600 & 2000 & $e$ in 2000\\
\hline
Saturn & $5:4$, major third & $+32$ & $+7$ & $-2$ & 0.0555\\
Jupiter & $6:5$, minor third & $+32$ & $+16$ & $+21$ & 0.0485\\
Mars & $3:2$, fifth & $-60$ & $-56$ & $-53$ & 0.0934\\
Earth & $16:15$, semitone & $+13$ & $+5$ & $+4$ & 0.0167\\
Venus & $25:24$, diesis & $-21$ & $-22$ & $-24$ & 0.0068\\
Mercury & $12:5$, octave + minor third & $-43$ & $-72$ & $-71$ & 0.2056\\
\hline
mean miss & & 33 & 30 & 29 & \\
chance of doing as well & & 0.22 & 0.15 & 0.14 & \\
\hline
\end{tabular}
\end{table}

The last line is the test of Section~\ref{sec:18resonance}: six ratios drawn at random between a unison
and a twelfth, each matched to the nearest of sixteen intervals of the just scale (the kinds Kepler drew
on), miss by $47$ cents on average and do as well as the planets about one time in six. By that test the
harmonies of the single planets are not established, in his figures or in ours. The Earth's semitone,
the ``mi, fa, mi'' that Kepler heard as misery and famine, and Saturn's major third are the close fits;
Mars and Mercury miss by half a semitone or more.

\emph{Did they fit better in his day?} No. The sky of 1600 held the same six planets on almost the same
orbits: the eccentricities have changed in the fourth decimal place, and the fastest mover, Saturn's
interval, drifts by about $2$ cents a century. The columns for 1600 and 2000 differ by $9$ cents for
Saturn and by $5$ or less for the rest. Where Kepler's own column differs from the orbit of 1600 the cause
is his data, and it cuts both ways. His figure for Saturn's slowest motion, $1'46''$ where the orbit of 1600 gives $1'48''$, cost him $25$
cents. His eccentricity for Mercury, $0.2095$ against a true $0.2056$, made Mercury fit $29$ cents better
than it should. For Mars, the Earth and Venus his speeds are within a quarter of a per cent of the modern
reconstruction, a tribute to Tycho Brahe's observations. (One caution about Saturn: Jupiter
disturbs its orbit so strongly that its eccentricity depends on the span averaged over. A fit to the years
1800 to 2050 alone gives a miss of $-13$ cents where the long formula gives $-2$.)

\emph{Why they cannot hold.} Nothing fixes an eccentricity at a musical value. The planets trade
eccentricity among themselves over tens and hundreds of thousands of years. The Earth's has ranged from
nearly $0$ to $0.068$ \cite{laskar2011}, which carries its ``semitone'' from a unison to within a
quarter-tone of a perfect fourth. On the present trend the Earth will sing a true $16:15$ around the year
3600, and Saturn's major third is passing through exactness about now. Period ratios are different: a
resonance holds them (Section~\ref{sec:18resonance}). The harmony Kepler was looking for exists where
there is something to hold it. One such place is the family of Jupiter's moons, which he was among the
first to study and for which he coined the word \emph{satellite}.

\emph{Two planets at a time.} Kepler set more store by the intervals between neighbours: the slowest motion
of the outer planet against the fastest of the inner (``diverging'') and the fastest of the outer against
the slowest of the inner (``converging''), ten intervals in all, such as $1:2$ for Saturn and Jupiter
converging and $2:3$ for Mars and the Earth. Here the test comes out differently. With the orbits of 2000
the ten intervals miss his consonances by $33$ cents on average, where random intervals miss the nearest
consonance by $57$; the chance of doing as well is about $0.03$. It is $0.015$ if the same six
eccentricities are dealt to the planets in a different order (all $720$ orders), and $0.03$ if the
eccentricities are drawn at random. As Kepler himself concluded, the converging intervals fit best ($23$
cents) and the diverging ones worst ($42$). Mars and the Earth converging make a fifth to within $2$
cents.

That would be a modest result in Kepler's favour, if it were the end of the evidence. It is not, because
planets were found after he died, and they are a test his scheme never had: Uranus (1781), Neptune
(1846), and the smaller bodies Ceres (1801) and Pluto (1930). As single voices Uranus and Neptune do
about as well as Jupiter and Venus ($9$ cents from a minor third and from a diesis), while Ceres and
Pluto miss by $47$ and $62$. But the four new intervals between neighbours, for Saturn with Uranus and
Uranus with Neptune, miss the nearest consonance by $86$, $55$, $154$ and $90$ cents, worse than chance.
With all fourteen intervals the mean miss is $51$ cents and the chance of doing as well is $0.29$. Kepler
also arrived at the angular speeds seen from the Sun only after trying and rejecting distances, periods
and linear speeds, which weakens any single test. So the fair verdict is that his two-planet harmonies
were a real feature of the six planets he knew, at about the level of one chance in thirty, and that the
planets he did not know do not continue it. (A further planet far beyond Neptune has been proposed since
2016. None had been found when this was written, so it cannot be tested.)

\paragraph{The second law in the atom.} Kepler's second law, equal areas in equal times, is the
conservation of angular momentum: a body of mass $m$ and angular momentum $L$ sweeps area at the rate
$L/(2m)$. Bohr's rule for the atom is $L=n\hbar$. So the electron in orbit $n$ sweeps area at
\[
  n\cdot\frac{\hbar}{2m_e}=n\times5.79\times10^{-5}\ \text{m}^{2}\text{/s},
\]
a whole number of times one fixed rate. (The Earth sweeps $2.2\times10^{15}$\,m$^{2}$/s.) Kepler's second
law is what Bohr quantized. In Sommerfeld's refinement of 1916 \cite{sommerfeld1916} the electron moves on
Kepler ellipses, and the ellipses allowed at level $n$ are those whose short and long axes are in the
ratio $k/n$, $k=1,\ldots,n$: the shapes of the orbits are fractions. That picture was replaced by quantum
mechanics ten years later, but the count of states and the symmetry of
Section~\ref{sec:18kepler} survived the change.

\paragraph{What else of Kepler's belongs here.} Five more pieces of his work touch the subjects of this
paper.

\emph{The golden ladder.} In his essay on the six-cornered snowflake \cite{kepler1611} Kepler observed
that the ratios of neighbouring Fibonacci numbers, $5:8$, $8:13$, $13:21$, come ever closer to the golden
section. Those ratios are the convergents of $[1;1,1,1,\ldots]$ (Chapter~\ref{ch:cf}): the golden movement of the
Suite, three centuries early.

\emph{Consonance and the totient.} In the first three books of \emph{Harmonices Mundi} Kepler derived the
consonances from the regular polygons that can be constructed with ruler and compass. The triangle,
square, pentagon, hexagon and octagon give the consonant ratios; the heptagon cannot be constructed and,
he argued, that is why $7$ gives no consonance \cite{gingerich1973}. Gauss and Wantzel later settled which polygons can be
constructed: the $n$-gon can exactly when $\varphi(n)$ is a power of $2$. That is the totient of
Chapter~\ref{ch:coprime}, the same function that weights the correlation matrix (\cref{prop:oq-gram}) and that appeared in
the ratio $n/\varphi(n)$ of Proposition~\ref{prop:18exact}. The polygons Kepler called knowable, with
$3,4,5,6,8,10,12,\ldots$ sides, are in modern terms among those with $\varphi(n)$ a power of $2$.

\emph{Tilings that never repeat.} Book II of \emph{Harmonices Mundi} lists the eleven tilings of the plane
by regular polygons, and draws a pattern of pentagons and stars that anticipates the non-repeating
tilings of Penrose (1974). In those tilings the two kinds of tile occur in the golden ratio, and since 1984 patterns of
this kind have been known in real materials, the quasicrystals \cite{Shechtman1984}. The equation of Section~\ref{sec:18butterfly} is the standard model
of an electron in a one-dimensional pattern of this kind.

\emph{Packing, and $\pi$.} The snowflake essay also asserts that no packing of equal spheres is denser
than the greengrocer's pile, which fills $\pi/\sqrt{18}=0.7405$ of space. That is how the atoms sit in
many metals. The claim was proved by Hales \cite{hales2005}.

\emph{The first inverse-square law.} In his optics of 1604 Kepler stated that the brightness of a light
falls off as the square of the distance. The forces of Section~\ref{sec:18kepler} follow the same law,
and it is the reason an orbit in the atom and an orbit in the sky are one problem.

Two tools of his are also at work in the instrument that accompanies this section. Kepler's equation,
$M=E-e\sin E$, gives the place of a planet on its ellipse at a given time; the instrument solves it for
every frame it draws. And his logarithms of 1624 are the natural measure of a musical interval: the cents
used above are a logarithm, and $\log_2(3/2)$ is one of the five numbers of the Suite.

%-------------------------------------------------------------------------------------------------------
\section{What is known and what is ours}\label{sec:18audit}

\begin{table}[ht]
\centering
\small
\begin{tabular}{p{0.34\textwidth}p{0.16\textwidth}p{0.42\textwidth}}
\hline
idea & class & notes\\
\hline
Kepler's law in the atom; the $2n^{2}$ shells from the Runge--Lenz symmetry & classical & Pauli, Fock.
The reading as a repair of the ``shared formula'' caveat of Chapters~\ref{ch:coprime} and~\ref{ch:connections} is ours.\\
Atoms with $n\approx1009$ and the fourth harmonic & specialist & Stepkin et al.; the arithmetic of the
harmonic is our check.\\
Orbital resonances, Kirkwood gaps, KAM & classical & Standard celestial mechanics.\\
The big-term test and Table~\ref{tab:18res} & new presentation & Continued fractions of period ratios are
an old device; the test against random ratios is our computation.\\
The sign mirror, Proposition~\ref{prop:18mirror} & new, as far as we found & A two-line consequence of
\cref{prop:th-inv}. May well be known for GCD matrices.\\
The one-prime chain and its mirror map & classical & The Kac--Murdock--Szeg\H{o} matrix and the Poisson
kernel.\\
Hofstadter's butterfly, Hall integers, Cantor spectrum & specialist & Hofstadter; Thouless et al.; Avila
and Jitomirskaya.\\
The butterfly read as a number's musical form; the instrument & new presentation & Our computation and
instrument.\\
The harmonic form, Proposition~\ref{prop:18harmonic} & new, as far as we found & One line from
\cref{prop:oq-gram}. The harmonic numbers make it worth stating.\\
The ratio of the two ends, Tables~\ref{tab:18ratio} and~\ref{tab:18smooth} & new computation & To
$m=3\times10^{7}$. The exact cases of Proposition~\ref{prop:18exact} are elementary.\\
Conjunction figures, Proposition~\ref{prop:18conj} & new presentation & The cycles are ancient and the
three-gap theorem is classical; reading the trigon as a convergent and its drift as the next term is ours.\\
Kepler's harmonies tested, Table~\ref{tab:18harm} & new computation & Others have re-examined them; the
tests against random ratios, the comparison of 1600 with 2000, and the test on the planets found later
are ours.\\
Constructible polygons and $\varphi(n)$; quasicrystals; sphere packing & classical & Gauss, Wantzel;
Penrose, Shechtman; Hales. Only the pointers to this paper are ours.\\
The second law in Bohr's atom; Sommerfeld's ellipses & classical & Bohr 1913, Sommerfeld 1916.\\
\hline
\end{tabular}
\end{table}

Still open, and the one item here that could become a theorem: does the ratio of $1/\lambda_{\min}$ to
$\lambda_{\max}$ tend to a limit, and if so to what? Propositions~\ref{prop:18mirror}
and~\ref{prop:18harmonic} reduce it to comparing the Perron eigenvalues of two explicit non-negative
matrices of the same kind, both built from the common multiples of pairs of rhythms. The computation of
Section~\ref{sec:18mirror} says the answer for the gcd matrix is close to $2.57$ at every size we can
reach, that the correlation matrix is converging on the same value from above, and that for any fixed
set of primes the answer is $1$.
